\documentclass[a4paper,12pt]{article}

\usepackage{fontspec}
\usepackage[russian]{babel}
\usepackage[a4paper,left=30mm,right=15mm,top=20mm,bottom=20mm]{geometry}
\usepackage{microtype}
\usepackage{setspace}
\usepackage{indentfirst}
\usepackage{amsmath}
\usepackage{booktabs}
\usepackage{siunitx}
\usepackage{xcolor}
\usepackage{tikz}
\usepackage{pgfplots}
\usepackage{pgfplotstable}
\usepackage{fancyhdr}
\usepackage{titlesec}
\usepackage{caption}
\usepackage{enumitem}
\usepackage{fancyvrb}
\usepackage[hidelinks]{hyperref}
\usepackage[russian,noabbrev]{cleveref}

\setmainfont{PTF55F.ttf}[
  BoldFont = PTF75F.ttf,
  ItalicFont = PTF56F.ttf,
  BoldItalicFont = PTF76F.ttf
]
\setsansfont{PTS55F.ttf}[
  BoldFont = PTS75F.ttf,
  ItalicFont = PTS56F.ttf,
  BoldItalicFont = PTS76F.ttf
]
\setmonofont{PTM55F.ttf}[
  BoldFont = PTM75F.ttf,
  Scale = 0.86
]
\emergencystretch=2em

\setstretch{1.28}
\setlength{\parindent}{1.25cm}
\setlength{\parskip}{0pt}

\usetikzlibrary{positioning,arrows.meta,calc,fit}
\pgfplotsset{compat=1.18}

% Okabe-Ito colorblind-friendly palette
\definecolor{oiBlue}{HTML}{0072B2}
\definecolor{oiTeal}{HTML}{009E73}
\definecolor{oiOrange}{HTML}{D55E00}
\definecolor{oiPurple}{HTML}{CC79A7}
\definecolor{oiSlate}{HTML}{5C6B73}
\definecolor{oiDark}{HTML}{1E293B}

% siunitx configuration
\sisetup{
  detect-all,
  output-decimal-marker = {.},
  group-separator = {\,},
  group-minimum-digits = 4,
  table-number-alignment = center,
  per-mode = symbol
}

% cleveref Russian labels
\crefname{figure}{рис.}{рис.}
\Crefname{figure}{Рис.}{Рис.}
\crefname{table}{табл.}{табл.}
\Crefname{table}{Табл.}{Табл.}
\crefname{section}{разд.}{разд.}
\Crefname{section}{Разд.}{Разд.}

% GOST / VKR section headings: PT Sans, centered, bold, uppercase
\titleformat{\section}[block]
  {\sffamily\fontsize{13}{15}\bfseries\centering}
  {\thesection.\ }
  {0pt}
  {\MakeUppercase}
\titlespacing*{\section}{0pt}{13pt plus 2pt minus 2pt}{7pt plus 1pt}

% GOST / VKR captions
% Figures: centered italic "Рис. N Название"
\captionsetup[figure]{
  labelsep=space,
  font={it,normalsize},
  labelfont={it},
  justification=centering,
  skip=6pt
}
\addto\captionsrussian{\renewcommand{\figurename}{Рис.}}

% Tables: "Таблица N" right-aligned, then centered title
\DeclareCaptionFormat{gosttable}{\hfill#1\par\centering#3\par}
\captionsetup[table]{
  format=gosttable,
  labelsep=none,
  font={normalsize},
  labelfont={normalfont},
  skip=6pt
}

% Compact lists with GOST indentation
\setlist[itemize]{leftmargin=1.8cm,itemsep=2pt,topsep=4pt,parsep=0pt}
\setlist[enumerate]{leftmargin=1.8cm,itemsep=2pt,topsep=4pt,parsep=0pt}

% Page numbering ONLY in bottom-right corner, no rules, no headers
\pagestyle{fancy}
\fancyhf{}
\fancyfoot[R]{\normalsize\thepage}
\renewcommand{\headrulewidth}{0pt}
\renewcommand{\footrulewidth}{0pt}

% Unified TikZ styles
\tikzset{
  sysbox/.style={
    rectangle,
    rounded corners=2pt,
    draw=black!75,
    line width=0.6pt,
    inner sep=5.5pt,
    align=center,
    font=\sffamily\footnotesize
  },
  agentbox/.style={sysbox, fill=blue!5, draw=oiBlue!80, line width=0.8pt},
  neutralbox/.style={sysbox, fill=gray!8},
  tealbox/.style={sysbox, fill=teal!6, draw=oiTeal!80},
  gatebox/.style={sysbox, fill=red!5, draw=oiOrange!85, line width=0.7pt},
  flowarrow/.style={-{Stealth[length=5pt]}, line width=0.65pt, draw=black!75},
  alertarrow/.style={-{Stealth[length=5pt]}, line width=0.7pt, draw=oiOrange!85},
  edgelabel/.style={font=\tiny\ttfamily, fill=white, inner sep=1.5pt, rounded corners=1pt}
}

\begin{document}

% ==================== TITLE BLOCK (GOST / ITMO) ====================
\begin{center}
  {\sffamily\small\MakeUppercase{Университет ИТМО}}\par
  \smallskip
  {\sffamily\fontsize{14}{18}\bfseries Ассистент сервисной службы <<ПолярТех>>:\\
  Graph~RAG и агент с инструментами ERP}\par
  \smallskip
  {\small
    Модель LLM: \texttt{qwen3:4b-instruct-2507-q4\_K\_M} (Ollama) \quad|\quad
    Эмбеддинги: \texttt{bge-m3}
  }
\end{center}

% ==================== 1. ЗАДАЧА ====================
\section{Задача}\label{sec:task}

Компания <<ПолярТех>> обслуживает коммерческое и промышленное холодильное оборудование (витрины, моноблоки, чиллеры). Диспетчер сервисной службы принимает заявки, содержащие модель оборудования и код ошибки панели управления (например, <<\texttt{AR-210}, код \texttt{E-17}>>). Задача ассистента --- оперативно определить отказавший узел, регламент ремонта и актуальную деталь с учётом снятия устаревших позиций с производства и цепочек замен (\texttt{REPLACED\_BY}), проверить остатки по \num{4} складам (Москва, Санкт-Петербург, Казань, Новосибирск) и цены в учётной системе (ERP), а затем подготовить сервисный заказ с обязательным подтверждением человеком.

Плоский векторный поиск (Flat~RAG) системно ошибается в двух классах сценариев:
\begin{enumerate}
  \item \textbf{Многошаговые замены (\texttt{dev-06}):} на вопросе <<Какую деталь заказать при ошибке \texttt{E-17} на \texttt{AR-210}?>> плоский RAG извлекает фрагмент руководства \texttt{manual\_AR-210.md} и отвечает снятой с производства деталью \texttt{P-1042}, не проходя по сервисному бюллетеню \texttt{bulletin\_SB-2025-03.md} к актуальной замене \texttt{P-1042B}.
  \item \textbf{Оперативные данные ERP (\texttt{dev-17}, \texttt{dev-20}):} на вопросе <<Сколько штук детали \texttt{P-1118B} есть на складе в Казани?>> плоский RAG получает точность \num{0.000}, так как динамические остатки и цены отсутствуют в статических документах и доступны только через API учётной системы.
\end{enumerate}

% ==================== 2. РЕШЕНИЕ ====================
\section{Архитектура решения и ReAct-конвейер}\label{sec:arch}

Система объединяет три взаимодействующих контура (\cref{fig:arch}), реализованных в соответствии с требованиями FR-1--FR-10:
\begin{itemize}
  \item \textbf{Извлечение и граф знаний (FR-1, FR-2):} модуль \texttt{extraction.py} извлекает сущности и типизированные связи по строгой JSON~Schema (\texttt{EXTRACTION\_SCHEMA}) с лимитом \num{4096} токенов и \num{4} синтетическими few-shot примерами без утечки эталона \texttt{gold\_dev.json}. Все коды заземляются на исходный чанк (\texttt{find\_codes}), названия нормализуются по словарю (\texttt{canonical\_name}), а для плановых регламентов без связей изолированные узлы \texttt{Procedure} не добавляются. Модуль \texttt{graph\_builder.py} строит ориентированный граф \texttt{nx.DiGraph} (\num{118} узлов, \num{279} рёбер), фиксируя источники на обоих концах ребра и разрешая взаимные конфликты \texttt{REPLACED\_BY} в пользу ребра с большим весом.
  \item \textbf{Сообщества и два режима поиска (FR-3--FR-5):} модуль \texttt{communities.py} разбивает граф методом Лувена (\texttt{seed=7}) на \num{8} сообществ. Для сообществ с линейками оборудования код детерминированно подсчитывает лидеров отказов и цепочки снятых деталей, а для тематических сообществ без линеек запрещает модели выдумывать линейки и рейтинги. Локальный поиск (\texttt{local\_search}) собирает пути длиной $\le 3$ между опорными узлами (или 2-hop окрестность при одном якоре), автоматически разворачивает цепочки \texttt{REPLACED\_BY} вместе со связями \texttt{SUPPLIED\_BY} и \texttt{IS\_A} и мгновенно возвращает отказ (\texttt{refused=True}) при отсутствии кода в графе. Глобальный поиск (\texttt{global\_search}) проверяет неизвестные коды до запуска цикла и выполняет Map-Reduce (\texttt{MAP\_SCHEMA}) по сводкам сообществ.
  \item \textbf{Инструменты и ReAct-агент (FR-6--FR-10):} модуль \texttt{tools.py} предоставляет \num{5} типизированных инструментов с валидацией аргументов (\num{4} склада ERP, формат кода детали). Модуль \texttt{agent.py} ведёт ReAct-цикл (\texttt{max\_steps=6}), нормализует коды, проверяет заземление деталей и чисел, выполняет до \num{3} попыток при \texttt{ErpUnavailable} и перехватывает \texttt{create\_service\_order} при \texttt{confirm=False} в \texttt{pending\_action}.
\end{itemize}

\begin{figure}[htbp]
  \centering
  \begin{tikzpicture}[node distance=6mm and 7mm, transform shape, scale=0.95]
    \node[neutralbox, text width=24mm] (disp) {\textbf{Диспетчер}\\Модель: \texttt{AR-210}\\Ошибка: \texttt{E-17}};

    \node[agentbox, right=of disp, text width=35mm] (agent) {%
      \textbf{ReAct-агент}\\
      Thought--Action цикл\\
      Лимит \texttt{max\_steps=6}\\
      Заземление кодов и чисел\\
      Дедупликация шагов
    };

    \node[neutralbox, right=11mm of agent, text width=56mm] (erp) {%
      \textbf{Учётная система (ERP)}\\
      \texttt{check\_stock} (\num{4} склада), \texttt{get\_price}\\
      До \num{3} попыток при \texttt{ErpUnavailable}
    };

    \node[tealbox, above=5mm of erp, text width=56mm] (graph) {%
      \textbf{Граф знаний (Graph~RAG)}\\
      \num{118} узлов, \num{279} рёбер, \num{8} сообществ\\
      \texttt{local\_search} (пути $\le 3$ + \texttt{REPLACED\_BY}) / Map-Reduce
    };

    \node[gatebox, below=5mm of erp, text width=56mm] (gate) {%
      \textbf{Шлюз безопасности (FR-9, FR-10)}\\
      \texttt{confirm=False}: перехват \texttt{create\_service\_order}\\
      Возврат \texttt{pending\_action} диспетчеру
    };

    \draw[flowarrow] (disp.east) -- (agent.west);
    \draw[flowarrow] (agent.east) -- (erp.west);
    \draw[flowarrow] (agent.east) -- ++(5mm,0) |- (graph.west);
    \draw[alertarrow] (agent.east) -- ++(5mm,0) |- (gate.west);
  \end{tikzpicture}
  \caption{Архитектура ассистента <<ПолярТех>>: взаимодействие ReAct-агента с графом знаний, ERP и шлюзом подтверждения человеком}\label{fig:arch}
\end{figure}

% ==================== 3. РЕЗУЛЬТАТЫ ====================
\section{Результаты и сравнительный бенчмарк}\label{sec:results}

Оценка выполнена на \num{24} контрольных вопросах \path{questions_dev.jsonl} (\texttt{temperature=0}, \texttt{seed=7}, чистый кеш \path{artifacts/llm_cache/}, модель \texttt{qwen3:4b-instruct-2507-q4\_K\_M}). Итоговые показатели четырёх режимов приведены в \cref{tab:bench} и визуализированы на \cref{fig:bench}.

\begin{table}[htbp]
  \centering
  \singlespacing
  \caption{Сравнение четырёх режимов работы ассистента на \num{24} вопросах \texttt{questions\_dev.jsonl}}\label{tab:bench}
  \scriptsize
  \setlength{\tabcolsep}{4pt}
  \begin{tabular}{@{} l S[table-format=1.3] S[table-format=1.3] S[table-format=1.3] S[table-format=1.3] @{}}
    \toprule
    \textbf{Метрика} & {\textbf{Baseline (Flat)}} & {\textbf{Local (Graph)}} & {\textbf{Global (Graph)}} & {\textbf{Agent (Граф+ERP)}} \\
    \midrule
    \textbf{Общая точность (Accuracy)} & 0.500 & \bfseries 0.792 & 0.417 & \bfseries 0.917 \\
    \quad Вопросы \texttt{local} (\num{5}) & 1.000 & \bfseries 1.000 & 0.800 & \bfseries 1.000 \\
    \quad Вопросы \texttt{multihop} (\num{7}) & 0.571 & \bfseries 1.000 & 0.143 & \bfseries 1.000 \\
    \quad Вопросы \texttt{global} (\num{4}) & 0.000 & \bfseries 1.000 & 0.500 & 0.750 \\
    \quad Вопросы \texttt{tool} (\num{5}) & 0.000 & 0.000 & 0.000 & \bfseries 0.800 \\
    \quad Вопросы \texttt{refusal} (\num{3}) & 1.000 & \bfseries 1.000 & \bfseries 1.000 & \bfseries 1.000 \\
    \midrule
    Выдуманные факты (Fabricated) & \bfseries 0.000 & {0.042$^*$} & 0.042 & 0.042 \\
    Полнота источников (Source Recall) & 0.856 & \bfseries 1.000 & 0.719 & \bfseries 1.000 \\
    Отклонённые вызовы инструментов & {---} & {---} & {---} & 0.032 \\
    Вызовов модели на вопрос & {1.00} & {\bfseries 0.88} & {7.67} & {4.83} \\
    Токенов на вопрос & {\num{1356.7}} & {\bfseries \num{1113.5}} & {\num{3887.9}} & {\num{5130.4}} \\
    \bottomrule
  \end{tabular}
\end{table}

\begin{figure}[htbp]
  \centering
  \begin{tikzpicture}
    \begin{axis}[
      name=accplot,
      ybar=1.2pt,
      bar width=4pt,
      width=0.58\textwidth,
      height=50mm,
      ymin=0, ymax=1.24,
      ylabel={Доля верных ответов},
      ylabel style={font=\scriptsize},
      title={\textbf{Точность по типам вопросов}},
      title style={font=\footnotesize},
      symbolic x coords={Local (5), Multihop (7), Global (4), Tool (5), Refusal (3), Все (24)},
      xtick=data,
      xticklabel style={font=\tiny, rotate=14, anchor=east},
      yticklabel style={font=\scriptsize},
      axis x line*=bottom,
      axis y line*=left,
      ymajorgrids=true,
      grid style={dashed, black!20},
      axis on top=false,
      enlarge x limits=0.11,
      legend style={at={(0.5,0.99)}, anchor=north, legend columns=4, font=\tiny, draw=black!20, fill=white}
    ]
      \addplot[fill=oiSlate!85, draw=none] table[x=category, y=baseline, col sep=comma]{report/benchmark_accuracy.csv};
      \addplot[fill=oiTeal!85, draw=none] table[x=category, y=local, col sep=comma]{report/benchmark_accuracy.csv};
      \addplot[fill=oiPurple!85, draw=none] table[x=category, y=global, col sep=comma]{report/benchmark_accuracy.csv};
      \addplot[fill=oiBlue!85, draw=none] table[x=category, y=agent, col sep=comma]{report/benchmark_accuracy.csv};
      \legend{Flat RAG, Local, Global, Agent}
    \end{axis}

    \begin{axis}[
      at={($(accplot.south east)+(12mm,0)$)},
      anchor=south west,
      ybar,
      bar width=12pt,
      width=0.38\textwidth,
      height=50mm,
      ymin=0, ymax=6200,
      ylabel={Токенов на вопрос},
      ylabel style={font=\scriptsize},
      title={\textbf{Стоимость режимов}},
      title style={font=\footnotesize},
      symbolic x coords={Local, Flat RAG, Global, Agent},
      xtick=data,
      xticklabel style={font=\tiny},
      yticklabel style={font=\scriptsize},
      axis x line*=bottom,
      axis y line*=left,
      ymajorgrids=true,
      grid style={dashed, black!20},
      axis on top=false,
      enlarge x limits=0.22,
      nodes near coords,
      nodes near coords style={font=\tiny\bfseries, color=oiDark},
      every node near coord/.append style={/pgf/number format/fixed, /pgf/number format/precision=0}
    ]
      \addplot[fill=oiBlue!75, draw=none] table[x=mode, y=tokens, col sep=comma]{report/benchmark_tokens.csv};
    \end{axis}
  \end{tikzpicture}
  \caption{Точность четырёх режимов по категориям вопросов и удельная стоимость в токенах на вопрос}\label{fig:bench}
\end{figure}

Звёздочкой в \cref{tab:bench} отмечено единственное ложное срабатывание метрики \texttt{fabricated} в режиме \texttt{local} на повторе <<\SI{2}{\text{шт.}}>> из формулировки вопроса \texttt{dev-21} (\cref{sec:errors}).

Качество извлечения на \num{16} эталонных документах \path{gold_dev.json} (\path{eval_extraction.json}): связи --- Precision~\num{1.000}, Recall~\num{1.000} (\num{161} из \num{161}); сущности --- Precision~\num{1.000}, Recall~\num{1.000} (\num{165} из \num{165}); отброшено проверками заземления (\texttt{dropped}): \num{0}.

Стоимость индексации (\path{index_stats.json}): \num{80} документов (\num{92} чанка), \num{118} узлов, \num{279} рёбер, \num{8} сообществ, \num{0} отброшенных элементов, \num{0} конфликтов. Затраты LLM: \num{100} вызовов (\num{20} попаданий в кеш), \num{228614} prompt-токенов и \num{40812} completion-токенов. Повторный прогон всех \num{5} оценок по закоммиченному кешу полностью воспроизводит числа \cref{tab:bench}.

% ==================== 4. ГДЕ ГРАФ ПОМОГАЕТ И ГДЕ УСТУПАЕТ ====================
\section{Где граф помогает и где уступает}\label{sec:compare}

\textbf{Пример~1 (преимущество графа и агента --- \texttt{dev-06}, \texttt{dev-17}, \texttt{dev-22}--\texttt{dev-24}).} На вопросе <<Ошибка \texttt{E-17} на \texttt{AR-210}. Какую деталь заказать, есть ли она на складе в Казани и сколько она стоит?>> плоский RAG останавливается на устаревшей детали \texttt{P-1042} и не имеет доступа к ERP. В графе знаний (\cref{fig:subgraph}) \texttt{local\_search} проходит путь длиной \num{3} ребра \texttt{AR-210 -> P-1042 -> пусковое реле <- E-17}, автоматически разворачивает цепочку \texttt{P-1042 -> REPLACED\_BY -> P-1042B -> SUPPLIED\_BY -> ХолодКомплект}, после чего агент запрашивает ERP и возвращает точный ответ (\SI{3}{\text{шт.}}, \SI{2630}{\text{руб.}}). На вопросах отказа (\texttt{dev-22}--\texttt{dev-24}) граф мгновенно возвращает \texttt{refused=True} за \num{0} вызовов LLM против \num{1357} токенов у плоского RAG.

\textbf{Пример~2 (где граф уступает по стоимости или точности --- \texttt{dev-01}--\texttt{dev-05}, \texttt{dev-14}).} На простых однохоповых вопросах \texttt{local} (\texttt{dev-01}--\texttt{dev-05}) плоский RAG достигает тех же \num{100}\% (\num{5}/\num{5}) за \num{1} вызов модели без затрат на предварительную графовую индексацию (\num{269426} токенов), тогда как агент тратит в среднем \num{4.83} вызова (\num{5130.4} токена). На обзорном вопросе \texttt{dev-14} (<<Какой узел чаще всего отказывает в линейке <<Гольфстрим>>?>>) маршрутизация через \texttt{summarize\_topic} (Map-Reduce) привела к путанице модели-лидера (\texttt{GS-510}, \num{4} отказа) с узлом-лидером (<<циркуляционный насос>>, \num{8} отказов), тогда как прямой \texttt{local\_search} по подграфу линейки ответил верно.

\begin{figure}[htbp]
  \centering
  \begin{tikzpicture}[node distance=9mm and 14mm]
    \node[agentbox, text width=25mm] (model) {\scriptsize Model:\\\textbf{\texttt{AR-210}}};
    \node[neutralbox, dashed, right=of model, text width=28mm] (pold) {\scriptsize Part (снята):\\\textbf{\texttt{P-1042}}};
    \node[tealbox, right=of pold, text width=28mm, line width=0.8pt] (pnew) {\scriptsize Part (замена):\\\textbf{\texttt{P-1042B}}};

    \node[gatebox, below=of model, text width=25mm] (err) {\scriptsize ErrorCode:\\\textbf{\texttt{E-17}}};
    \node[sysbox, fill=orange!6, draw=oiOrange!70, below=of pold, text width=28mm] (comp) {\scriptsize Component:\\\textbf{пусковое реле}};
    \node[sysbox, fill=purple!5, draw=oiPurple!80, below=of pnew, text width=28mm] (supp) {\scriptsize Supplier:\\\textbf{ХолодКомплект}};

    \node[sysbox, fill=cyan!5, draw=oiBlue!60, below=7mm of err, text width=25mm] (proc) {\scriptsize Procedure:\\\textbf{\texttt{SP-104}}};

    \draw[flowarrow] (model) -- node[edgelabel]{HAS\_PART} (pold);
    \draw[alertarrow] (pold) -- node[edgelabel, text=oiOrange]{\textbf{REPLACED\_BY}} (pnew);
    \draw[flowarrow] (pold) -- node[edgelabel]{IS\_A} (comp);
    \draw[flowarrow] (err) -- node[edgelabel]{INDICATES} (comp);
    \draw[flowarrow] (pnew) -- node[edgelabel]{SUPPLIED\_BY} (supp);
    \draw[flowarrow] (err) -- node[edgelabel]{FIXED\_BY} (proc);
  \end{tikzpicture}
  \caption{Фрагмент графа знаний вокруг \texttt{AR-210}: путь длиной \num{3} ребра и разворачивание цепочки \texttt{REPLACED\_BY}}\label{fig:subgraph}
\end{figure}

% ==================== 5. АНАЛИЗ ОШИБОК ====================
\section{Анализ ошибок системы и оптимизация}\label{sec:errors}

В ходе итеративной отладки были разобраны ошибки на всех трёх этапах конвейера:
\begin{enumerate}
  \item \textbf{Этап извлечения (\texttt{extraction.py}):} при лимите \texttt{max\_tokens=1024} длинные руководства обрывались на середине JSON (Recall связей \num{0.329}, \num{50} рёбер). После увеличения лимита до \num{4096} Recall вырос до \num{0.590}, но на \num{5} из \num{12} руководств (\texttt{AR-310}, \texttt{AR-320}, \texttt{BZ-210}, \texttt{GS-510}, \texttt{GS-520}) модель имитировала 1-пунктовый пример из промпта и извлекала только первую деталь. Добавление второго пункта в синтетический пример и явного правила извлечения всех строк довело качество до \num{1.000}/\num{1.000} и \num{279} рёбер (\cref{tab:extraction}).
  \item \textbf{Этап сообществ и поиска (\texttt{communities.py}, \texttt{search.py}):} при общем промпте сводок тематическое сообщество компрессора (без \texttt{ProductLine}) выдумывало <<узел-лидер \texttt{E-11}>> и перебивало реальные линейки в \texttt{global\_search}. Разделение промптов для сообществ с линейками и без них устранило шум, однако на \texttt{dev-14} (\texttt{summarize\_topic}) модель всё ещё путает модель-лидер \texttt{GS-510} с узлом-лидером. В режиме \texttt{local} на \texttt{dev-21} повтор <<\SI{2}{\text{шт.}}>> из формулировки вопроса даёт \num{1} ложное срабатывание метрики \texttt{fabricated} (\num{0.042}), поскольку \texttt{evaluation.py} сверяет числа с единицами только с результатами инструментов.
  \item \textbf{Этап агента и инструментов (\texttt{agent.py}):} на \texttt{dev-21} агент верно определил \texttt{P-2054C}, проверил остаток (\SI{6}{\text{шт.}}) и цену (\SI{6750}{\text{руб.}}), но вместо вызова \texttt{create\_service\_order} сгенерировал текстовый ответ с расчётной суммой и <<\SI{2}{\text{шт.}}>> --- встроенный фильтр \texttt{\_find\_discrepancies} дважды перехватил неподтверждённые числа и безопасно перевёл ответ в отказ (\texttt{refused=True}). На \texttt{dev-12} при верном ответе (\texttt{P-1118B}, <<ТермоПрибор>>) модель исказила код регламента (\texttt{SP-205} вместо \texttt{SP-105}), что составило единственный реальный выдуманный факт агента (\num{1}/\num{24} = \num{0.042}).
\end{enumerate}

\begin{table}[htbp]
  \centering
  \singlespacing
  \caption{Эволюция качества извлечения на \num{16} эталонных документах \texttt{gold\_dev.json} и размере графа}\label{tab:extraction}
  \scriptsize
  \setlength{\tabcolsep}{4.5pt}
  \begin{tabular}{@{} l S[table-format=1.3] S[table-format=1.3] S[table-format=3.0] S[table-format=2.0] @{}}
    \toprule
    \textbf{Конфигурация \texttt{extraction.py}} & {\textbf{Precision}} & {\textbf{Recall}} & {\textbf{Рёбер в графе}} & {\textbf{Отброшено (\texttt{dropped})}} \\
    \midrule
    Лимит \texttt{max\_tokens=1024} & 1.000 & 0.329 & 50 & 12 \\
    \texttt{max\_tokens=4096}, \num{1} пример & 1.000 & 0.590 & 125 & 38 \\
    \textbf{\texttt{max\_tokens=4096}, \num{4} синтетических примера} & \bfseries 1.000 & \bfseries 1.000 & \bfseries 279 & \bfseries 0 \\
    \bottomrule
  \end{tabular}
\end{table}

% ==================== 6. ГРАНИЦЫ СИСТЕМЫ ====================
\section{Границы системы и безопасность (FR-9, FR-10)}\label{sec:safety}

Контур безопасности (\cref{fig:safety}) исключает несанкционированные изменения в ERP и галлюцинации в аргументах:
\begin{itemize}
  \item \textbf{Подтверждение человеком (FR-9):} при \texttt{confirm=False} (по умолчанию) вызов \texttt{create\_service\_order} перехватывается до обращения к ERP и не добавляется в список выполненных шагов; агент возвращает \texttt{needs\_confirmation=True} и структуру \texttt{pending\_action} с параметрами заказа. Реальное создание заказа выполняется только при явном запуске с флагом \texttt{--confirm}.
  \item \textbf{Заземление аргументов (FR-7, FR-10):} номер детали \texttt{part\_id} должен присутствовать в вопросе пользователя или в результатах \texttt{search\_knowledge}; количество \texttt{qty} --- в вопросе; склад проверяется по закрытому списку из \num{4} городов (Москва, Санкт-Петербург, Казань, Новосибирск). При сбоях \texttt{ErpUnavailable} выполняются до \num{3} попыток; ответы чётко различают <<нет данных>> и <<остаток \num{0}>>.
\end{itemize}

\begin{figure}[htbp]
  \centering
  \begin{tikzpicture}[node distance=6mm and 6mm, transform shape, scale=0.92]
    \node[neutralbox, text width=23mm] (call) {\textbf{Вызов модели}\\(\texttt{ToolCall})};
    \node[neutralbox, right=of call, text width=28mm] (params) {\textbf{1.~Параметры}\\Склад (\num{4} города)\\Формат \texttt{Part~ID}};
    \node[neutralbox, right=of params, text width=29mm] (ground) {\textbf{2.~Заземление}\\\texttt{part\_id} в вопросе\\или в графе?};
    \node[agentbox, right=of ground, text width=28mm] (optype) {\textbf{3.~Тип операции}\\Чтение ERP или\\создание заказа?};

    \node[tealbox, above=7mm of optype, text width=28mm] (read) {\textbf{Чтение ERP~API}\\\texttt{/stock}, \texttt{/parts}};
    \node[gatebox, below=10mm of optype, text width=30mm] (stop) {\textbf{\texttt{confirm=False}: STOP}\\\texttt{pending\_action} (FR-9)};
    \node[gatebox, below=10mm of $(params.south)!0.5!(ground.south)$, text width=58mm] (rej) {\textbf{ОТКЛОНЕНИЕ:} возврат ошибки модели (FR-7)};

    \draw[flowarrow] (call) -- (params);
    \draw[flowarrow] (params) -- (ground);
    \draw[flowarrow] (ground) -- (optype);
    \draw[flowarrow, draw=oiTeal!85] (optype) -- (read);
    \draw[alertarrow] (optype) -- (stop);
    \draw[alertarrow] (params.south) -- (params.south |- rej.north);
    \draw[alertarrow] (ground.south) -- (ground.south |- rej.north);
  \end{tikzpicture}
  \caption{Блок-схема барьеров безопасности и шлюза Human-in-the-Loop}\label{fig:safety}
\end{figure}

Ниже приведена реальная трассировка агента из \path{artifacts/eval_agent.json} (\texttt{dev-17}) и схема срабатывания шлюза подтверждения:
\begin{singlespace}
\begin{Verbatim}[fontsize=\scriptsize,frame=single,rulecolor=\color{black!30},framesep=4pt,samepage=true]
Вопрос: «Ошибка E-17 на AR-210. Какую деталь заказать, есть ли она в Казани и сколько стоит?»
  1. search_knowledge({'question': 'Ошибка E-17 на AR-210'})
     -> пусковое реле; P-1042 заменена на P-1042B (ХолодКомплект), регламент SP-104
  2. check_stock({'part_id': 'P-1042B', 'warehouse': 'Казань'})
     -> {'part_id': 'P-1042B', 'warehouse': 'Казань', 'qty': 3, 'discontinued': False}
  3. get_price({'part_id': 'P-1042B'})
     -> {'part_id': 'P-1042B', 'name': 'пусковое реле', 'price_rub': 2630, 'discontinued': False}
Ответ: «Заменить деталь P-1042B (пусковое реле). Склад: Казань, остаток: 3 шт., цена: 2630 руб.»

При запросе заказа (confirm=False):
  [Safety Gate] create_service_order(part_id='P-2054C', warehouse='Москва', qty=2, model='BZ-120')
  -> needs_confirmation=True, pending_action={'tool': 'create_service_order', 'args': {...}}
\end{Verbatim}
\end{singlespace}

% ==================== 7. СЛЕДУЮЩИЙ ШАГ ====================
\section{Следующий шаг развития системы}\label{sec:next}

В качестве проверяемой гипотезы предлагается структурированный маршрутизатор заказов в \texttt{agent.py} в сочетании с типизацией сущностей в \texttt{MAP\_SCHEMA} для \texttt{global\_search}. Если в запросе присутствует намерение оформления заказа и все \num{4} аргумента уже подтверждены, агент подсказывает модели прямой вызов \texttt{create\_service\_order} (\texttt{dev-21}); а добавление поля \texttt{entity\_type} (\texttt{component} / \texttt{model} / \texttt{part}) в \texttt{MAP\_SCHEMA} исключает смешение модели-лидера и узла-лидера на шаге \texttt{reduce} (\texttt{dev-14}). Критерий проверки: рост точности агента с \num{0.917} (\num{22}/\num{24}) до \num{1.000} (\num{24}/\num{24}) и снижение \texttt{fabricated\_rate} агента до \num{0.000}.

% ==================== 8. ВКЛАД УЧАСТНИКОВ ====================
\section{Вклад участников}\label{sec:contrib}

Проект выполнен индивидуально: проектирование промптов и заземления (\texttt{extraction.py}, \texttt{graph\_builder.py}); кластеризация Лувена, детерминированная статистика сообществ и оба режима поиска (\texttt{communities.py}, \texttt{search.py}); ReAct-агент, типизированные инструменты ERP и шлюз подтверждения (\texttt{tools.py}, \texttt{agent.py}); \num{41} модульный тест в \path{tests/unit/} (\num{183} теста всего с контрактными и инфраструктурными); развёртывание локального Ollama, полный бенчмаркинг \num{4} режимов, анализ ошибок и вёрстка итогового отчёта.

% ==================== 9. ИСПОЛЬЗОВАНИЕ ИИ ====================
\section{Использование ИИ-инструментов}\label{sec:ai}

В процессе разработки использовались ИИ-ассистенты кодирования (Antigravity / Gemini): (1)~генерация модульных тестов граничных случаев (циклы \texttt{REPLACED\_BY}, сбои \texttt{ErpUnavailable}, перехват \texttt{confirm=False}, отказ на неизвестных кодах); (2)~диагностика усечения JSON при \texttt{max\_tokens=1024} и неполного извлечения списков руководств, рефакторинг промптов без утечки эталона; (3)~подготовка LaTeX-вёрстки отчёта с векторными схемами TikZ и графиками PGFPlots. Все архитектурные решения, код \num{6} модулей, результаты прогонов, метрики и целостность защищённой инфраструктуры проверены автором лично.

\end{document}
